\section{System Architecture}
\label{sec:architecture}

Two rules constrain the design. First, every decision procedure (an ordering, an identity, a hash,
a predicate) has one implementation, compiled for both the host and the device. Second, no thread
waits for another: every shared structure is lock-free, and where threads must agree, one
compare-exchange decides and the losers adopt the winner's result.

Algorithm~\ref{alg:evolution} gives the evolution loop. A match task for each new state runs the
join of Algorithm~\ref{alg:matching} and creates a rewrite task for each accepted match. A rewrite
task applies its match, claims the child's canonical identity, records the event and its causal
and branchial relations, forwards the surviving matches (Section~\ref{sec:forwarding}), and
creates a match task for the child if it was first to claim the child's identity. Under a per-state
cap or sampling, the choice among a state's matches waits until the state's matching is complete
(Section~\ref{sec:sampling}), and under quotient exploration only one state per class gets a match
task (Section~\ref{sec:quotient}).

\begin{algorithm}[htbp]
\caption{Task-Based Multiway Evolution}
\label{alg:evolution}
\begin{algorithmic}[1]
\Procedure{Evolve}{$L$, $S_0$, depth $D$}
    \ForAll{$s_0 \in S_0$}
        \State publish $s_0$; enqueue \textsc{MatchTask}($s_0$)
    \EndFor
    \While{tasks pend at a depth $\le D$}
        \State $\tau \gets$ \textsc{Dequeue}()
        \If{$\tau = \textsc{MatchTask}(s)$}
            \State $M \gets$ inherited $\cup$ \textsc{FindMatches}($L$, $s$)
            \ForAll{deduplicated $m \in M$}
                \State enqueue \textsc{RewriteTask}($s$, $m$)
            \EndFor
        \ElsIf{$\tau = \textsc{RewriteTask}(s, m)$}
            \State apply $m$, creating raw child $c$
            \State canonicalize $c$; claim its identity
            \State record event, causal and branchial edges
            \State forward $s$'s surviving matches to $c$
            \If{claim won and $\mathrm{depth}(c) < D$}
                \State enqueue \textsc{MatchTask}($c$)
            \EndIf
        \EndIf
    \EndWhile
\EndProcedure
\end{algorithmic}
\end{algorithm}

\subsection{Shared Host and Device Implementation}
\label{sec:shared-core}

The CPU engine and the CUDA port have separate schedulers and share their semantics. Each decision
procedure they share is one function compiled for host and device from a single source: the join
that enumerates matches, the edge-signature rule, the rewrite rule that determines a produced
edge's vertices, the event identity, the content-ordered state key, the individualization--refinement
canonical hash, and the quotient-causal dynamic program and its per-instance replay. There are
eleven such shared functions; the other three are the matcher's task core, the sampling draw and
the slot rule. The differential tests between host and device therefore cover what differs between
them: scheduling, memory placement and queueing.

\subsection{Unified Hypergraph Storage}
\label{sec:storage}

All edges are stored in one shared pool, and each state is a bitset over that pool
(Figure~\ref{fig:storage}), so a child that differs from its parent by a few edges stores only the
difference.

\begin{figure}[htbp]
\centering
\begin{tikzpicture}[
    box/.style={rectangle, draw, minimum width=2cm, minimum height=0.6cm},
    label/.style={font=\small}
]
% Edge Pool
\node[box, minimum width=6cm] (pool) at (0,0) {Edge Pool (SegmentedArray)};
\node[label, above=0.1cm of pool] {$e_0, e_1, e_2, \ldots, e_n$};

% States
\node[box] (s1) at (-2.5,-1.5) {State $s_1$};
\node[box] (s2) at (0,-1.5) {State $s_2$};
\node[box] (s3) at (2.5,-1.5) {State $s_3$};

% Arrows
\draw[->] (s1.north) -- node[left, label] {bitset} (s1.north |- pool.south);
\draw[->] (s2.north) -- node[left, label] {bitset} (s2.north |- pool.south);
\draw[->] (s3.north) -- node[right, label] {bitset} (s3.north |- pool.south);

% Legend
\node[label, below=0.3cm of s2] {States are views (SparseBitset) over shared edges};
\end{tikzpicture}
\caption{Unified edge storage with state views}
\label{fig:storage}
\end{figure}

The bitset is divided into chunks, and a child shares its parent's chunks until it changes one, at
which point that chunk is copied. A rewrite that changes a few edges allocates a few chunks, and
shared chunks are never modified, so readers of the parent are unaffected by the child.
Table~\ref{tab:t3-deheap} gives the arena bytes of each corpus workload. Appendix~\ref{app:data-structures}
specifies the lock-free structures, their per-operation costs, the memory model, and the behavior
at configured limits.

\subsection{Allocation}
\label{sec:allocation}

Evolution allocates from per-worker arenas. Each worker bumps a pointer into blocks it owns, so
allocation during matching and rewriting needs no synchronization, and memory is freed in bulk by
resetting the arena. The global allocator is called only to obtain those blocks:
Table~\ref{tab:t3-deheap} reports $455$ to $475$ calls for an entire run, which on the largest
workloads is $0.007$ calls per event.

\subsection{Static Rule Analysis}
\label{sec:static-analysis}
\label{sec:rule-analysis}

When a rule is added, the engine computes properties that follow from its patterns alone and uses
them to choose a strategy. For each rule it records the difference in edge count between the two
sides, which says whether states grow; the number of new vertices the right-hand side introduces;
and the left-hand side's arities. Treating the left-hand side as a conjunctive query, it computes
whether the query is acyclic and an integral edge cover $\rho$, which bounds the matches in a state
of $N$ edges by $N^{\rho}$. The bound is tight when the query is acyclic, and is used only as an
upper bound otherwise.

Each predicate is sound in one direction. The branching test is an over-approximation: a negative
answer means no two matches can ever share a consumed edge, and a positive answer means only that
this has not been ruled out. The engine acts only on the negative answer: the branchial relation
is then empty for every initial state, and the engine returns it empty without computing it. No
analysis here decides termination, finiteness of the reachable set, or confluence, which are
undecidable for rewriting systems of this kind.

\subsection{Determinism Contract}
\label{sec:determinism}

The state set, the event set and the causal and branchial relations are a deterministic function of
the rules, the initial state, the step count and the options. They do not depend on the number of
worker threads, the order in which tasks run, or whether quotient exploration is enabled. State and
event identifiers are allocated in arrival order and vary between runs; the observables are defined
over canonical identities. A determinism gate checks this across thread counts and seeds: it
fingerprints the expanded state set, the multiset of (input, output, rule) transitions, and the
causal and branchial relations separately, and reports which component differs and under which
configuration. The components are compared separately because the causal relation can differ by a
single edge while the others agree.

The count caps \textsc{MaxStatesPerStep} and \textsc{MaxSuccessorStatesPerParent} are outside the
contract. They admit states until a counter reaches the cap, so which states are admitted depends
on arrival order and varies with the worker count. Making them schedule-independent would require
holding every candidate at a depth until the depth is complete, which is a barrier. A capped run
guarantees the count at every worker count, not which states make it up.

\subsection{Depth Semantics and Termination}
\label{sec:depth}

The step budget applies to the \emph{shortest} path from an initial state to a canonical state, not
to the path along which the state was first reached. Under concurrency a state can be reached first
along a long path and later along a short one, so a budget on arrival depth would include or exclude
states according to which path won the race, and the explored set would vary with the thread count.

The engine therefore \emph{relaxes} depths: when a state is reached along a path shorter than its
recorded depth, the recorded depth is lowered and the change is propagated to its children, and
from them onward. This terminates because each accepted change lowers a depth, which is bounded
below by zero. A state is matched at most once however often its depth is lowered, because the
claim that admits a state to expansion is separate from its depth.

Relaxation races the registration of a new child. Registration publishes the child and then reads
the parent's depth; relaxation lowers the depth and then scans the child list. Sequentially
consistent fences on both sides rule out the case where the scan misses the child and the read
misses the lowered depth, so a child always sees its parent's current minimum depth. A state past
the budget stays on the frontier without being claimed, so a shorter path found later can still
bring it within the budget. A state is claimed against the same bound its match task is admitted
under, so a claimed state is never left unmatched.

Work budgets are taken by compare-and-swap. An increment followed by a rollback would briefly
publish a count above the limit, and a thread reading the count at that moment would prune work
that another schedule keeps.

Termination is detected without a barrier: a depth is complete when it has no live work and the
depth below it is complete. This suffices because a task at depth $d$ only creates work at depths
above $d$, so once depth $d-1$ is complete nothing new can arrive at $d$.

\subsection{Transitive Reduction of the Causal Graph}
\label{sec:tr}

The causal graph has an edge for every dependency between events, including those implied by
longer paths. Its transitive reduction removes an edge $p \to c$ whenever a path
$p \rightsquigarrow c$ of length greater than one exists; on an acyclic graph the reduction is
unique.

Goranci et al.~\cite{goranci2025dynamic} give the first fully dynamic algorithms for the transitive
reduction of a general directed graph under insertions and deletions, in $\BigO(m + n\log n)$
amortized update time, and $\BigO(m + n^{1.585})$ worst case with fast matrix multiplication. This
engine's case is simpler in two ways. First, causal edges are only inserted, so an edge dropped as implied
never needs to be reconsidered. Second, edges arrive in topological order, because an event's producers
exist before the event, so when an edge $(p,c)$ arrives the closure below $p$ is final. The
procedure below is the specialization of the problem to this case.

\paragraph{Online maintenance.}
The reduction is kept up to date after every insertion, because a session
(Section~\ref{sec:sessions}) can be queried between steps, and because building the full relation
first would allocate the transitive closure that the reduction avoids.

\paragraph{Procedure (Algorithm~\ref{alg:tr}).}
Let $e = (p, c)$ be a new causal edge: event $p$ produced an edge that event $c$ consumed. If $c$ is
already reachable from $p$ through retained edges, $e$ is implied and dropped; otherwise it is
retained. A retained edge never needs to be removed later. A consumer's incoming edges are added by
its own worker in descending producer order, and an ancestor's id is always smaller than its
descendant's. A retained $(p', c)$ could become implied only through a path
$p' \rightarrow \cdots \rightarrow x \rightarrow c$. If $(x, c)$ arrived first, the query for
$(p', c)$ saw the path and dropped the edge; if $(x, c)$ arrived later, then $x < p'$ by the
descending order, while reachability from $p'$ requires $p' < x$, a contradiction. Hence deciding each
edge once on arrival gives the exact reduction. The descending order also makes queries cheap,
because the closure from nearer producers is built before farther ones are tested, and it needs no
sorting, because events are numbered in creation order.

\begin{algorithm}[htbp]
\caption{Online Transitive Reduction (insertion-only)}
\label{alg:tr}
\begin{algorithmic}[1]
\Procedure{ConsumeEdges}{event $c$, producers $P$}
    \ForAll{$p \in P$ in descending event id}
        \If{$c$ reachable from $p$ via retained edges}
            \State drop $(p, c)$ as implied
        \Else
            \State retain $(p, c)$
        \EndIf
    \EndFor
\EndProcedure
\end{algorithmic}
\end{algorithm}

\paragraph{Order independence.}
The retained edge set must depend only on the causal DAG, not on the order in which concurrent
workers insert edges; the determinism gate checks this across thread counts and seeds. Full capture
meets it through the arrival order above: each consumer's edges are added by its own worker, in
descending producer order, after its producers' edges are complete, and under that order the
decision on arrival is exact however consumers interleave. The quotient reconstruction adds edges
between canonical event ids, which are not ordered this way, so it stores the relation unreduced
and reduces it when it is read; a finite DAG has one transitive reduction, so the result depends
only on the stored edges.

Only the causal graph is reduced. States, events and the branchial relation are not.

